% ============================================================
% Section 4 — Query-Level Evidence Credit (v2, 2026-09-16)
% 主方法 = 查询级信用 (HANDOFF §25); 广播覆盖项降为对照臂 (4.4)。
% 超参: M=3, β_q=1.0, 查询封顶 16 词; 对照臂 β=0.3; λ_s=0.015, token 计价 1e-3, KL 0.02/0.5 → 附录 C
% v1 (2026-08-28 至 09-03, 广播覆盖项为主方法) 原文以注释保留在文末。
% ============================================================

\section{Query-Level Evidence Credit}
\label{sec:method}

\subsection{Trajectory return}
\label{sec:method-reward}

Each rollout $\tau$ receives a scalar trajectory return
\begin{equation}
R(\tau) \;=\; R_{\mathrm{ans}}(\tau)
\;-\; \lambda_{s}\, N_{\mathrm{search}}(\tau)
\;-\; P_{\mathrm{aux}}(\tau),
\label{eq:reward}
\end{equation}
where $R_{\mathrm{ans}}$ is the terminal answer reward (alias-aware exact
match against the gold answer set), $N_{\mathrm{search}}$ counts search
calls, and $P_{\mathrm{aux}}$ collects format, turn-limit, and free-token
penalties; \emph{free tokens} are the tokens a rollout emits outside
protocol actions (plan lines, search queries, answers). The estimator
computes one leave-one-out advantage $A_i$ per rollout from this return
(Eq.~\ref{eq:rloo}) and broadcasts it to the tokens of every turn of
$\tau_i$.

\subsection{Evidence gain of a query}
\label{sec:method-gain}

Let $\mathcal{G}$ be the question's annotated set of gold supporting
documents, matched by document title, and let $D_t$ be the documents
retrieved by the query $q_t$ issued at search turn $t$. The
\emph{evidence gain} of $q_t$ is the share of $\mathcal{G}$ that the
rollout reaches for the first time at turn $t$:
\begin{equation}
g(q_t) \;=\; \frac{\bigl|\,(D_t \cap \mathcal{G}) \setminus
\textstyle\bigcup_{t'<t} D_{t'}\,\bigr|}{|\mathcal{G}|}.
\label{eq:gain}
\end{equation}
Each gold document is credited once, at its first retrieval;
re-retrieving it earns nothing. The \emph{coverage} of a rollout,
\begin{equation}
C(\tau) \;=\; \sum_t g(q_t) \;=\; \frac{\bigl|\{\, d \in \mathcal{G} :
\tau \text{ retrieves } d \,\}\bigr|}{|\mathcal{G}|},
\label{eq:coverage}
\end{equation}
is the fraction of $\mathcal{G}$ it reaches over the whole trajectory.

\subsection{Query-level credit}
\label{sec:method-credit}

At each search turn $t$ of rollout $\tau_i$, let $s_t$ be the state---the
context preceding the query. We sample $M$ additional queries
$q'_1,\ldots,q'_M \sim \pi_\theta(\cdot \mid s_t)$, retrieve for each
against the same environment, and compute their evidence gains under
Eq.~\ref{eq:gain} with the same prior-hit set. The credit of the
executed query is
\begin{equation}
c_{i,t} \;=\; \beta_q \Bigl( g(q_t) - \frac{1}{M}\sum_{m=1}^{M}
g(q'_m) \Bigr).
\label{eq:credit}
\end{equation}
The sampled queries serve only as a baseline: they are not appended to
the trajectory, receive no gradient, and are discarded after retrieval.
The credit is added to the advantage of the tokens of turn $t$,
\begin{equation}
A_{i,t} \;=\; A_i + c_{i,t},
\label{eq:turn-adv}
\end{equation}
and answer turns keep $A_i$. The bracket in Eq.~\ref{eq:credit} is a
Monte-Carlo estimate of the advantage of $q_t$ over the policy's own
action distribution at $s_t$, with the baseline formed from samples
drawn at the visited state as in the vine estimator of
\citet{schulman2015trpo}. It is computed from the gold annotations and
the retriever; no value function is involved. Retrieval uses the first
16 words of a query; tokens beyond the cap do not reach the retriever,
so a query cannot raise its credit by appending terms. We use $M=3$ and
$\beta_q=1$ (Appendix~C).

In an answer-tied group, $A_i$ carries only the cost terms of
Eq.~\ref{eq:cost-ordering}. The credit $c_{i,t}$ is non-zero wherever
the executed query and the policy's alternatives at the same state
reach different shares of $\mathcal{G}$, which holds at 35--50\% of the
search states of the trained policy in the closed pool
(\S\ref{sec:mech-broadcast}). The tie is broken at the query step; the
answer step is left to the outcome reward.

\subsection{Trajectory-level coverage: the contrast arm}
\label{sec:method-effect}

The alternative placement adds coverage to the trajectory return,
\begin{equation}
R^{\mathrm{cov}}(\tau) \;=\; R(\tau) \;+\;
\mathbf{1}\!\left[\mathrm{ans}(\tau)\right]\, \beta\, C(\tau),
\label{eq:covreward}
\end{equation}
gated on the rollout committing to an answer: without the gate, an
abstaining rollout collects coverage credit, and in ungated runs
abstention rates rose with coverage reward. Its contribution to the
leave-one-out advantage is $\beta\,\tfrac{K}{K-1}\bigl(C_i -
\bar{C}\bigr)$ (Eq.~\ref{eq:cov-advantage}): a within-group ordering of
whole trajectories by coverage, broadcast to every turn. We train this
arm with $\beta=0.3$ as the controlled contrast; \S\ref{sec:mech-broadcast}
measures what the broadcast does at the level of turns.

\subsection{Stabilizing long-horizon optimization}
\label{sec:method-kl}

Trajectories mix two token populations: protocol actions and free
tokens. We regularize the two channels separately: a weak KL
coefficient on action tokens and a strong KL coefficient on free
tokens, both toward the SFT reference policy, which rarely emits free
tokens. Under a uniform coefficient, free tokens grow geometrically
until training collapses, and interventions on the gradient
channel---token pricing at several rates, full and asymmetric gradient
masking---leave the growth rate unchanged (six runs, five
interventions; Fig.~\ref{fig:ftk}, Appendix~B). With split
coefficients, free-token counts stay flat over the full schedule.
Replacing the reference policy with a mid-training checkpoint that has
a high free-token rate re-ignites inflation within seven episodes (mean
free tokens per turn $0.7 \rightarrow 7.2$); an otherwise-identical run
that keeps the SFT reference policy does not.

% ---- KL 图(用户 2026-09-03: 先放占位符) ----
% 成图: paper_assets/fig_ftk_forensics.{pdf,png}(生成器 fig_ftk_forensics.py)
% 入稿前: 图例 "two-tier KL anchor (ours)"、"re-anchored ref" 改成 reference
% 措辞; 然后把下面的 \fbox 占位换成
% \includegraphics[width=\columnwidth]{../paper_assets/fig_ftk_forensics.pdf}
\begin{figure}[t]
\centering
\fbox{\parbox{0.95\columnwidth}{\centering\vspace{2em}
[Figure placeholder: free tokens per turn (log scale) over training
episodes---uniform KL with no intervention, token cost, full and
asymmetric gradient masking, replaced reference policy; two-tier KL
toward the SFT reference policy.]\vspace{2em}}}
\caption{Free tokens per turn over training under a uniform KL
coefficient with five interventions, versus the two-tier KL toward the
SFT reference policy. [Placeholder.]}
\label{fig:ftk}
\end{figure}

\subsection{Training procedure}
\label{sec:method-algo}

Each training step proceeds as follows.
\begin{enumerate}\setlength\itemsep{1pt}
\item Sample $K$ rollouts per question from the current policy under the
  search protocol.
\item Parse actions and execute searches against the environment. At
  each search turn, sample $M$ queries from the same state, retrieve
  for each, and record first-time gold-title hits for the executed query
  and for each sample.
\item Compute trajectory returns by Eq.~\ref{eq:reward} and query
  credits by Eq.~\ref{eq:credit}.
\item Compute leave-one-out advantages within each question group,
  broadcast each trajectory's advantage to its tokens, and add the
  credit of each search turn to that turn's tokens (Eq.~\ref{eq:turn-adv}).
\item Apply the policy-gradient update with channel-specific KL
  regularization; a small auxiliary value head serves as a variance
  monitor (Appendix~C).
\end{enumerate}
Hyperparameters, infrastructure, and guard criteria are given in
Appendix~C.

% ---------------------------------------------------------------
% v1 原文(广播覆盖项为主方法), 供回溯:
% \section{Coverage-Augmented Trajectory Ranking}
% 4.1 Reward definition: R(τ) = R_ans + 1[ans(τ)] β C(τ) − λ_s N_search − P_aux;
%   "Each gold document is credited at most once, at its first retrieval
%   (\emph{first-hit}); re-retrieving an already-credited document earns
%   nothing." "\textbf{Completion gating}: the coverage credit is paid only if
%   the rollout commits to an answer. Without the gate, an abstaining rollout
%   collects coverage credit; in ungated runs, abstention rates rose with
%   coverage reward. \textbf{Determinism}: we compute coverage as the set
%   overlap between retrieved evidence and the training annotations; the
%   signal is deterministic and used only during training."
% 4.2 The term is a trajectory-ranking signal: "Coverage credit is emitted
%   incrementally, at the turn where each first hit occurs. With discounting
%   near one, its trajectory sum equals β C(τ) up to negligible correction, and
%   the estimator computes a single leave-one-out advantage per trajectory---so
%   the term acts on \emph{trajectory returns and their within-group ranking}.
%   Step-level process supervision assigns credit per step; this term assigns
%   credit per trajectory. ... what remains is a within-group ordering by
%   coverage. Uniformly failing rollouts are thereby ranked by how much of the
%   required evidence they reached, displacing the cost-driven ordering of §3."
% 更早原句(清扫前): "so the term rewards evidence \emph{progress}, not
%   retrieval volume." / "Two design choices matter. ... the expected signature
%   of this arbitrage. Gating ties the process credit to the behavior the task
%   actually demands." / "There is no learned scorer, no reference plan to align
%   against, and no model in the reward loop; ... nothing about the mechanism
%   is needed at test time." / "we make no claim that individual turns are
%   separately credited, and we do not use a process reward model."
% ---------------------------------------------------------------
